% Copyright 2026 Sreeram Anil
% SPDX-License-Identifier: CC-BY-4.0
% =============================================================================
%  Benchmark methodology, metrics and validation
% =============================================================================
\chapter{Benchmark methodology, metrics and validation}
\label{ch:metrics}

Chapter~\ref{ch:datasets} defined the data. This chapter defines what is measured on them, how it
is measured, and --- the part that is usually left out --- how we established that the measuring
instrument itself is correct. Every number quoted anywhere in this report comes from one of four
executables (\code{bench\_cell}, \code{bench\_pack}, \code{bench\_imu}, \code{bench\_soh}) driven by
\code{scripts/run\_all.sh}, whose command lines are reproduced in Appendix~\ref{app:build}. The
outputs are CSV files under \code{results/final}; the tables and figures in this report are
generated from them by \code{scripts/make\_tables.py} and \code{scripts/make\_figures.py} and are
never edited by hand.

\section{The harness}

\subsection{Registry and the estimator interface}

Estimators are not enumerated in the benchmark. Each family registers itself in its own translation
unit through a function \code{register\_<family>(EstimatorList\&)}, and
\code{make\_all\_estimators()} calls all of them in a fixed order, so the benchmark has no
compile-time knowledge of any particular method; adding one is a matter of one
\code{add\_cell\_estimator<MyFilter<EcmModel>>} line (\S\ref{sec:lib-newest}).

The interface an estimator must satisfy is the abstract class \code{estkit::Estimator}
(\code{battery/ecm\_model.hpp}): \code{reset(const EstimatorConfig\&)},
\code{step(Real i, Real v, Real T)}, \code{soc()}, and optionally \code{soc\_std()},
\code{voltage\_pred()}, \code{capacity\_Ah()}, \code{soc\_lower()}/\code{soc\_upper()}, plus
\code{state\_bytes()} and \code{name()}/\code{group()}. Offline methods derive instead from
\code{BatchEstimator} and implement \code{run(i, v, T, n, soc\_out, v\_pred\_out)} over the whole
record. Optional outputs default to \code{NaN} or $-1$, and the metric code skips them, so a method
that has no covariance is simply not scored on covariance.

\subsection{The timing convention}
\label{sec:me-timing-convention}

The single most important convention in the harness is what happens inside one call to
\code{step}. Following Plett~\cite{plett2004ekf3}, with $\vu_k = [\,i^{\text{m}}_k,\ T^{\text{m}}_k\,]^\T$
and $y_k = v^{\text{m}}_k$:

\begin{algorithm}[H]
\caption{\code{CellEstimator::step($i^{\text{m}}_k, v^{\text{m}}_k, T^{\text{m}}_k$)} --- one sampling period}
\begin{algorithmic}[1]
\Require filter state $(\xplus_{k-1},\Pplus_{k-1})$, stored $\vu_{k-1}$, new measurements at $k$
\State $\vu_k \gets [\,i^{\text{m}}_k,\ T^{\text{m}}_k\,]^\T$
\If{$k>0$} \State \textbf{time update:} \code{predict}$(\vu_{k-1})$ \Comment{uses the \emph{previous} input}
\EndIf
\State $\yhat_k \gets$ \code{predict\_measurement}$(\vu_k)$ \Comment{$h(\xminus_k,\vu_k)$, recorded for the residual metric}
\State \textbf{measurement update:} \code{update}$(y_k, \vu_k)$
\State store $\vu_{k-1}\gets\vu_k$
\Ensure $(\xplus_k, \Pplus_k)$; \code{soc()} returns $\xplus_k[\soc]$
\end{algorithmic}
\end{algorithm}

Two details matter. First, the time update uses $\vu_{k-1}$, not $\vu_k$: the plant is integrated
with a zero-order hold on the current over $[t_{k-1},t_k)$, so the input that drove the transition
into sample $k$ is the one measured at $k-1$. Using $\vu_k$ instead would introduce a systematic
half-sample lead and an artificial correlation between the process noise and the innovation.
Second, \code{predict\_measurement} is called \emph{before} \code{update}, so the recorded
$\yhat_k$ is the \emph{prior} prediction $h(\xminus_k,\vu_k)$ and the voltage residual metric of
\S\ref{sec:me-residual} is a genuine one-step-ahead prediction error and not a fitted residual.
At $k=0$ there is no time update; the filter is initialised by \code{init}$(\vx_0,\mP_0)$ inside
\code{reset} and the first measurement is used immediately.

\subsection{Seeds, repetitions and outputs}

The cell-level and pack benchmarks are run with three seeds ($s=1,2,3$). Because the seed drives
the load generator as well as the sensor noise (\S\ref{sec:ds-gust}), three seeds means three
different flights; the tables report the mean over seeds and the harness also records the standard
deviation. The single-seed profiles (\code{fixed\_wing}, \code{hover\_hold},
\code{ground\_charge}) and the float32 build use $s=1$ only.

Each (estimator, profile, scenario, seed) combination is run $R=2$ times and the \emph{minimum}
wall-clock time is kept (\S\ref{sec:me-timing}). Only the first repetition records the optional
outputs (covariance standard deviation, interval bounds, capacity), so the second repetition is a
slightly leaner loop; taking the minimum therefore reports the warm-cache, no-instrumentation cost,
which is the number an implementer cares about. The estimate trajectories themselves are identical
across repetitions --- every estimator is re-\code{reset} with the same configuration and the same
seed --- which is itself a determinism check.

Time series are written for seed~1 only and decimated by a factor of ten
(\code{--ts-decimate}, default 10): \SI{1}{\hertz} records for the cell and pack benchmarks,
\SI{10}{\hertz} for the attitude benchmark. Decimation is applied by \emph{sub-sampling}, not
averaging, so the plotted traces are exact samples of the full-rate trajectory; all metrics are
computed at full rate before decimation.

\section{Cell-level metrics}

Let $\soc_k$ be the true state of charge, $\hat\soc_k$ the estimate, $n$ the number of samples and
\begin{equation}
e_k \;=\; \hat\soc_k - \soc_k, \qquad k = 0,\dots,n-1
\label{eq:me-err}
\end{equation}
the signed SOC error. SOC is dimensionless; all SOC metrics are reported in \emph{percentage
points} in the tables (a value of \num{0.36} means \SI{0.36}{\percent} of full capacity, i.e.\
\SI{18}{\milli\ampere\hour} on a \SI{5}{\ampere\hour} cell). If any $e_k$ is not finite the run is
marked divergent and that sample is replaced by $e_k = 1$ so that the remaining statistics stay
computable.

\subsection{Accuracy}

\begin{align}
\text{RMSE} &= \sqrt{\frac{1}{n}\sum_{k=0}^{n-1} e_k^2}, \label{eq:me-rmse}\\
\text{RMSE}_{\text{ss}} &= \sqrt{\frac{1}{n-\lfloor n/2\rfloor}\sum_{k=\lfloor n/2\rfloor}^{n-1} e_k^2},
 \label{eq:me-rmsess}\\
\text{MAE} &= \frac{1}{n}\sum_{k=0}^{n-1} \abs{e_k}, \qquad
e_{\max} = \max_k \abs{e_k}, \qquad
e_{\text{final}} = e_{n-1}. \label{eq:me-mae}
\end{align}

RMSE and $\text{RMSE}_{\text{ss}}$ answer different questions and the report uses both. RMSE over
the whole run is dominated by the transient whenever there is one: in \code{init\_error} a filter
that converges perfectly in \SI{10}{\second} still carries the \num{0.35} error for those seconds.
$\text{RMSE}_{\text{ss}}$, over the second half of the run --- which for the eVTOL mission begins at
$t=\SI{1005}{\second}$, mid-cruise, and includes the landing hover --- is the metric for
\emph{tracking} quality once the transient is over. It is the primary ranking metric of this report,
because a BMS spends almost all of its life in steady state and because the number that matters
operationally, the reserve at the end of the flight, is set by the error late in the mission.
$e_{\max}$ and $e_{\text{final}}$ are reported alongside because the mean-square metrics hide
exactly the behaviour a certification argument cares about: a filter with an excellent RMSE and a
\SI{5}{\percent} excursion during the landing hover is not acceptable, and $e_{\text{final}}$ is the
error at the moment the pilot decides whether to go around.

\subsection{Convergence}

\begin{definition}[Convergence time]\label{def:me-conv}
$t_{\text{conv}}$ is the earliest time $t=k\dt$ such that $\abs{e_j} < 0.02$ for all
$j$ in the window $[k,\ k + \lceil 60/\dt\rceil - 1]$, i.e.\ the start of the first
\SI{60}{\second} interval throughout which the SOC error stays inside \SI{2}{\percent}. If no such
window exists, $t_{\text{conv}} = -1$ and the tables print ``--''.
\end{definition}

Both the threshold and the dwell are deliberate: a threshold-crossing definition without a dwell
would reward a filter that oscillates through the band, and \SI{60}{\second} at \SI{10}{\hertz} is
600 consecutive samples --- long enough to exclude a chance crossing, short enough to resolve
differences between fast filters on a \SI{33.5}{\minute} mission. $t_{\text{conv}} = 0$ is the
common case in the nominal scenario, not a defect.

\subsection{Divergence}

\begin{definition}[Divergence]\label{def:me-div}
A run is \emph{divergent} if any $\hat\soc_k$ or $\hat v_k$ is non-finite, or if
$\abs{e_k} > 0.30$ for any $k \ge \lfloor n/2\rfloor$.
\end{definition}

Restricting the \SI{30}{\percent} test to the second half is what makes the criterion usable in
the \code{init\_error} scenario, where the error \emph{starts} at \num{0.35} by construction: a
filter that removes it is not divergent, one that still carries it half-way through the mission is.
The criterion is intentionally coarse. It answers ``did this run produce a number a supervisory
function could act on?'', not ``is this filter stable?''; stability claims are made in the
individual chapters from the structure of the algorithm, not from the benchmark.

\subsection{The voltage residual}
\label{sec:me-residual}

\begin{equation}
\text{RMSE}_v \;=\; \sqrt{\frac{1}{\abs{\mathcal{K}}}\sum_{k\in\mathcal{K}} \bigl(\hat v_k - v_k\bigr)^2},
\qquad \mathcal{K} = \{k : \hat v_k \text{ finite}\},
\label{eq:me-rmsev}
\end{equation}
where $\hat v_k = h(\xminus_k,\vu_k)$ is the prior voltage prediction recorded in
\S\ref{sec:me-timing-convention} and $v_k$ is the \emph{true} terminal voltage --- not the measured
one. This is a deliberate departure from the usual practice of reporting the innovation
$y_k - \hat v_k$. Comparing against the truth removes the sensor noise from the metric, so
$\text{RMSE}_v$ measures how well the estimator's internal model reproduces the cell, and a value
near zero means the model \emph{and} the state are right. Comparing against the measurement would
floor the metric at $\sigma_v$ and would reward a filter that has learned to reproduce the noise.
The figures, in contrast, plot the innovation $y_k - \hat v_k$, because that is the signal the
filter actually sees.

$\text{RMSE}_v$ is the single most useful diagnostic in the report, because it separates the two
ways an SOC estimate can be wrong. A filter with a large SOC error and a \emph{small} voltage
residual has found a state that explains the data --- the error is an observability or bias
problem, and no amount of better filtering will fix it. A filter with a large voltage residual has
a model that cannot explain the data at all. Compare, on the nominal eVTOL mission, Coulomb
counting ($\text{RMSE}_v = \SI{3.41}{\milli\volt}$, model fine, state drifting) with the same
method on \code{aged\_cell} ($\SI{112}{\milli\volt}$, model wrong).

\subsection{Capacity and interval metrics}

Estimators that report a capacity are scored on
\begin{equation}
\Delta Q \;=\; \frac{1}{\abs{\mathcal{L}}}\sum_{k\in\mathcal{L}}\bigl(\hat Q_k - Q_k\bigr),
\qquad \mathcal{L} = \{\,k \ge n - \lfloor n/10 \rfloor\,\},
\label{eq:me-cap}
\end{equation}
the mean \emph{signed} capacity error over the last \SI{10}{\percent} of the run, in
\si{\ampere\hour}. The signed mean, rather than an RMS, is used because a capacity estimate that
has converged to the wrong value is a bias, and its sign tells the reader whether the estimator is
optimistic (reports more capacity than exists --- the dangerous direction) or conservative.

Interval and set-membership observers, which report $[\hat\soc^-_k, \hat\soc^+_k]$ rather than a
distribution, are scored on
\begin{equation}
\text{viol} = \frac{1}{n}\,\bigl|\{k : \soc_k \notin [\hat\soc^-_k,\hat\soc^+_k]\}\bigr|,
\qquad
\bar w = \frac{1}{n}\sum_k \bigl(\hat\soc^+_k - \hat\soc^-_k\bigr).
\label{eq:me-bounds}
\end{equation}
These two must always be read together: a bound violation fraction of zero is trivial to achieve
with $\bar w = 2$, and the value of a guaranteed bound lies entirely in how tight it is.

\section{Pack-level metrics}

With $N_c = 12$ cells, per-cell errors $e_{c,k} = \hat\soc_{c,k} - \soc_{c,k}$:
\begin{align}
\text{RMSE}_{\text{cells}} &= \sqrt{\frac{1}{n N_c}\sum_{k}\sum_{c} e_{c,k}^2},
\qquad e_{\max} = \max_{c,k}\abs{e_{c,k}}, \label{eq:me-pack-cells}\\
\text{RMSE}_{\min} &= \sqrt{\frac{1}{n}\sum_k\Bigl(\min_c\hat\soc_{c,k} - \min_c\soc_{c,k}\Bigr)^2},
 \label{eq:me-pack-min}\\
\text{RMSE}_{\text{mean}} &= \sqrt{\frac{1}{n}\sum_k\Bigl(\tfrac{1}{N_c}\textstyle\sum_c\hat\soc_{c,k}
 - \tfrac{1}{N_c}\sum_c\soc_{c,k}\Bigr)^2}, \label{eq:me-pack-mean}\\
\text{RMSE}_{\max} &= \sqrt{\frac{1}{n}\sum_k\Bigl(\max_c\hat\soc_{c,k} - \max_c\soc_{c,k}\Bigr)^2}.
 \label{eq:me-pack-max}
\end{align}
The three aggregate metrics score the three operational definitions
\eqref{eq:ds-socmin}--\eqref{eq:ds-socmean}, and they are not redundant: the
$\min$ and $\max$ are order statistics, so an estimator that is unbiased per cell can still be
biased in $\text{RMSE}_{\min}$ (the minimum of $N_c$ noisy estimates is biased low by roughly
$1.6\,\sigma$ for $N_c=12$ Gaussian errors of standard deviation $\sigma$). This bias is a
\emph{property of the reporting definition}, not of the estimator, and one of the results of
Part~XI is how much each pack architecture manages to suppress it.

Pack runs additionally record \code{messages\_per\_step} --- the number of scalars a distributed
implementation would have to exchange between nodes per sample --- which is the currency in which
consensus and federated architectures pay for their accuracy, and \code{bytes} for the whole
$12$-cell estimator.

\section{Attitude metrics}

Attitude error is not an angle difference; it is a rotation. With $\hat q_k$ the estimate and $q_k$
the truth (both unit quaternions, body to navigation), the natural metric on $SO(3)$ is the
geodesic distance
\begin{equation}
\alpha_k \;=\; \bigl\|\log\!\bigl(q_k^{-1}\otimes\hat q_k\bigr)\bigr\|
 \;=\; 2\,\bigl|\!\arctan\!\bigl(\|\bm{\epsilon}\|,\ \epsilon_0\bigr)\bigr|,
\qquad q_k^{-1}\otimes\hat q_k = [\epsilon_0,\ \bm{\epsilon}^\T]^\T,
\label{eq:ds-imu-angdist}
\end{equation}
reported in degrees, with the sign of the error quaternion normalised so that $\epsilon_0 \ge 0$
(the double cover of $SO(3)$ by the unit quaternions means $q$ and $-q$ are the same rotation).
The benchmark reports $\text{RMS}(\alpha)$, its steady-state value over the second half,
$\max_k\alpha_k$, and a convergence time defined exactly as in Definition~\ref{def:me-conv} but
with a threshold of \SI{2}{\degree} and a dwell of \SI{30}{\second}. A run is divergent if
$\alpha_k > \SI{45}{\degree}$ in the second half.

Per-axis errors are reported as well, from the wrapped Euler differences
$\Delta\phi,\Delta\theta,\Delta\psi \in (-180^\circ,180^\circ]$, because they say \emph{which}
degree of freedom failed: a filter that loses yaw during a magnetic disturbance and one that loses
roll during a turn have the same $\alpha$ but entirely different causes. Euler differences are
only meaningful away from the pitch singularity; the dataset never exceeds \SI{12}{\degree} of
pitch, so this is safe here and would not be in an aerobatic dataset.

Gyro-bias accuracy is
\begin{equation}
\text{RMSE}_b = \sqrt{\frac{1}{n}\sum_k \bigl\|\hat{\bm b}^g_k - {\bm b}^g_k\bigr\|^2}
\label{eq:me-bias}
\end{equation}
in \si{\degree\per\second}. Estimators that do not model a bias report $\hat{\bm b}^g = \bm 0$, so
their $\text{RMSE}_b$ is the norm of the true bias --- a meaningful baseline rather than a missing
entry.

\section{Multi-flight state-of-health metrics}

For flight $f = 1,\dots,F$ with $F=120$, let $\hat Q_f$ be the capacity estimate averaged over the
last \SI{10}{\percent} of the flight phase and $Q_f$ the true capacity:
\begin{equation}
\text{RMSE}_Q = \sqrt{\frac{1}{F-10}\sum_{f=11}^{F}\bigl(\hat Q_f - Q_f\bigr)^2},
\qquad
\text{MAE}_Q = \frac{1}{F-10}\sum_{f=11}^{F}\bigl|\hat Q_f - Q_f\bigr|,
\label{eq:me-soh}
\end{equation}
both in \si{\ampere\hour}, with the first ten flights excluded so that the initial convergence of
the parameter estimate does not swamp a metric that is meant to describe tracking. The final
capacity error $\hat Q_F - Q_F$ and the mean per-flight SOC RMSE are reported alongside, the latter
to expose the common failure in which a joint estimator buys capacity accuracy with SOC accuracy.

\section{Execution time}
\label{sec:me-timing}

\subsection{What is measured}

The timed region (\code{benchmarks/bench\_cell.cpp}) brackets \code{reset()} and the complete sample loop:
for each $k$, \code{step(i,v,T)} followed by the accessor calls \code{soc()} and
\code{voltage\_pred()}. It is measured with \code{std::chrono::steady\_clock}, divided by the number
of samples, and minimised over $R$ repetitions:
\begin{equation}
t_{\text{step}} \;=\; \min_{r=1..R}\ \frac{T_r}{n}, \qquad R = 2 .
\label{eq:me-nsstep}
\end{equation}
The minimum, rather than the mean, is the right statistic for a deterministic computation on a
shared machine: every source of variation (scheduling, interrupts, frequency scaling, cache
eviction by another process) can only \emph{add} time, so the minimum is the best available
estimator of the true cost. The measurement therefore excludes nothing that the estimator does and
includes: the \code{reset()} call amortised over $n$ samples (negligible at $n=\num{20100}$ except
for estimators that solve a Riccati equation at reset, where it is the point), the virtual dispatch
through \code{Estimator*}, and the two accessor calls. It excludes dataset generation, metric
computation and all file output.

\subsection{Machine and build}

All timings in this report were taken on a single core of an x86-64 machine (an Intel Xeon-class
server core with a nominal base clock of \SI{2.1}{\giga\hertz}), with GCC at
\code{-O2}, C++17, \code{-march=native} \emph{disabled} (the default \code{ESTKIT\_NATIVE=OFF}), and
\code{Real = double}. Disabling \code{-march=native} is deliberate: it keeps the generated code to
the SSE2 baseline, which is far closer to what a scalar embedded target will do than an
AVX-512-vectorised build would be, and it makes the numbers reproducible across machines of the
same family. No multi-threading is used anywhere in the library.

\subsection{Scaling to a Cortex-M7 --- and why you should not trust the result}

The natural question for an airborne implementer is what these numbers become on the target. A
defensible order-of-magnitude translation rests on three explicit assumptions, and it is worth
being blunt that each of them can be wrong by a factor of two.

\begin{assumption}[Clock ratio]\label{as:me-clock}
The reference machine runs at $f_{\text{host}} \approx \SI{2.1}{\giga\hertz}$ base (higher under
turbo); a Cortex-M7 in a flight-grade BMS runs at
$f_{\text{M7}} = \SIrange{200}{480}{\mega\hertz}$. The clock ratio alone is thus
\numrange{4}{12}.
\end{assumption}

\begin{assumption}[Microarchitecture ratio]\label{as:me-ipc}
The host is a wide out-of-order core with several floating-point pipelines, out-of-order load/store
and large caches; the Cortex-M7 is a six-stage, dual-issue, in-order core with an FPv5 unit that
performs one double-precision multiply--accumulate over several cycles and runs from
tightly-coupled memory or flash with a cache~\cite{arm2021cortexm7}. For the scalar,
dependency-chained, small-matrix code in this library --- which has almost no instruction-level
parallelism to extract --- a further factor of \numrange{3}{6} is a reasonable expectation.
\end{assumption}

\begin{assumption}[No memory-system surprise]\label{as:me-mem}
The estimator's working set (a few kilobytes, Table~\ref{tab:me-mem}) fits in TCM or cache, so no
flash-wait penalty applies. This holds for the filters in this report and fails immediately for the
particle filters, whose state runs to tens of kilobytes.
\end{assumption}

Under Assumptions~\ref{as:me-clock}--\ref{as:me-mem} the translation factor is
\begin{equation}
t_{\text{M7}} \;\approx\; \kappa\, t_{\text{step}}, \qquad \kappa \in [\,20,\ 90\,],
\quad\text{central estimate } \kappa \approx 30 .
\label{eq:me-scaling}
\end{equation}
So the EKF's \SIrange{340}{630}{\nano\second} becomes roughly \SIrange{10}{20}{\micro\second} on a
\SI{200}{\mega\hertz} M7, i.e.\ \SIrange{0.01}{0.02}{\percent} duty at \SI{10}{\hertz} per cell, or
\SIrange{0.12}{0.24}{\percent} for a twelve-cell module --- comfortable. A bootstrap particle
filter at \SI{40}{\micro\second} per step becomes \SI{1.2}{\milli\second}, which is \SI{1.2}{\percent}
duty for one cell and \SI{14}{\percent} for a module: still feasible, but no longer free. Those
conclusions --- ``comfortable'' and ``no longer free'' --- are robust to a factor of two in
$\kappa$, and that is all \eqref{eq:me-scaling} is used for in this report. \textbf{No timing claim
in this document is a substitute for measurement on the target}, and a Cortex-M4, whose FPU is
single-precision only, will execute the \code{double} build in software emulation and be
\emph{ten to twenty times} slower again --- which is the practical argument for the float32 build of
\S\ref{sec:me-float}, not a numerical one.

\section{Memory footprint}

Memory is reported as \code{sizeof(Filter)} through \code{Estimator::state\_bytes()}: the complete
estimator object, including the \emph{model copy} each templated filter carries. That is what a
static allocation costs, but it needs one caveat, because for the battery problem it is dominated by
something that is not the filter.

\begin{table}[H]\centering\small
\caption{Where the bytes go for the ECM cell model (double precision, x86-64).}
\label{tab:me-mem}
\begin{tabular}{@{}lrl@{}}\toprule
object & \code{sizeof} [B] & note \\ \midrule
\code{OcvTable<241>} & 3872 & $2\times241$ doubles: the OCV table and its segment slopes \\
\code{CellParams} & 176 & 22 scalar parameters \\
\code{EcmModel} & 4104 & the table, the parameters, $\dt$, noise scalars, $q$-floors \\
\code{Ekf<EcmModel>} & 4336 & the model copy plus $\vx,\mP,\veps,\mS,\mK$: $232$\,B of filter state \\
\code{CoulombCounting} & 4184 & the model copy plus four scalars \\
\bottomrule\end{tabular}
\end{table}

For $n_x = 4$, $n_y = 1$ the filter's own state is $n_x + n_x^2 + n_y^2 + n_x n_y = 25$ reals
$= \SI{200}{\byte}$; every other byte is the model. A production BMS would hold one copy of the OCV
table in flash and share it across all twelve cells, turning a 12-cell decentralised EKF from
\SI{52}{\kilo\byte} into about \SI{6}{\kilo\byte}, so the reader should mentally subtract
\SI{3.9}{\kilo\byte} per instance when sizing a real target. Where the \emph{filter} state genuinely
dominates --- the particle, ensemble and Gaussian-sum methods, the moving-horizon estimators, the
fixed-lag smoothers --- the number needs no correction.

\section{The single-precision build}
\label{sec:me-float}

The library compiles unchanged with \code{-DESTKIT\_FLOAT=ON}, which sets
\code{ESTKIT\_REAL=float} and makes \code{estkit::Real} a 32-bit IEEE~754 binary32
value~\cite{ieee754_2019}. The whole benchmark is run a second time in that configuration
(\code{results/final/float}) and the results chapter reports the ratio
$\text{RMSE}_{\text{ss}}(\text{float32})/\text{RMSE}_{\text{ss}}(\text{float64})$ per estimator.

This is not a curiosity. A Cortex-M4 and the smaller M7 variants have a single-precision FPU only;
running \code{double} there means software emulation at a cost of an order of magnitude, so a
32-bit build is often the only viable one. The question the float32 build answers is which
estimators \emph{survive} it. The classical analysis is due to Verhaegen and Van
Dooren~\cite{verhaegen1986}, who showed that the conventional covariance recursion has an error
growth governed by the condition number $\kappa(\mP)$, and that the square-root and
$U$--$D$ factorised forms effectively halve the number of digits lost, because they propagate a
factor whose condition number is $\sqrt{\kappa(\mP)}$. The practical rule that follows is
\begin{equation}
\kappa(\mP) \;\lesssim\; \frac{1}{\varepsilon_{\text{mach}}}
\quad\text{for the conventional form}, \qquad
\kappa(\mP) \;\lesssim\; \frac{1}{\varepsilon_{\text{mach}}^2}
\quad\text{for the factored forms},
\label{eq:me-vvd}
\end{equation}
with $\varepsilon_{\text{mach}} = 2^{-24} \approx \num{6e-8}$ in binary32 (so
$1/\varepsilon_{\text{mach}} \approx \num{1.7e7}$) and $2^{-53}\approx\num{1.1e-16}$ in binary64.
The peak condition number of $\mP$ along the eVTOL mission is of order $10^6$ for the cell model,
uncomfortably close to the binary32 limit and utterly safe in binary64; this is precisely the regime
in which the square-root EKF (Chapter~\ref{ch:srekf}), the $U$--$D$ filter
(Chapter~\ref{ch:udekf}) and the square-root UKF (Chapter~\ref{ch:srukf}) earn their extra
arithmetic, and the float32 table is where that claim is tested rather than asserted.

Two implementation details change with \code{Real}. The finite-difference step used by the
numeric Jacobian fallback (\code{core/model.hpp}) switches from $10^{-6}$ to
$2\times10^{-3}$, following the $\varepsilon^{1/3}$ rule for central
differences~\cite{golubvanloan2013}; and the unit-test tolerances are relaxed to
$\max(\text{tol}, 2\times10^{-4})$. Nothing else in the library is conditioned on the precision.

\section{Reproducibility}

The benchmark is deterministic in the strong sense: a given \code{results/final} directory can be
reproduced bit-for-bit from the source tree and \code{scripts/run\_all.sh}. Three properties are
enforced by construction. \emph{Deterministic pseudo-randomness}: all randomness --- load gusts,
sensor noise, cell-parameter spreads, particle-filter resampling, ensemble initialisation --- comes
from \code{estkit::Rng}, an implementation of xoshiro256$^{**}$ with \code{splitmix64}
seeding~\cite{blackman2021xoshiro} that uses only 64-bit integer operations with defined wraparound,
so the bit stream is identical on every platform and compiler; Gaussian variates come from the
Marsaglia polar method~\cite{marsaglia1964polar}. No use is made of \code{std::random\_device}, of
the standard library's distribution objects (whose algorithms are implementation-defined), or of the
system clock. \emph{Stream separation}: the load, the sensors and the cell spreads draw from
separate \code{Rng} instances with different seed derivations (\S\ref{sec:ds-gust}), and stochastic
estimators construct their own generator from \code{EstimatorConfig::seed}, so a particle filter's
resampling draws never disturb the dataset. \emph{No hidden state}: \code{reset()} reconstructs the
model and the filter from scratch, which is what makes the repetition loop of \S\ref{sec:me-timing}
a determinism check as well as a timing measurement.

The one thing that is \emph{not} bit-reproducible across machines is the timing column; the
\code{real\_type} column in every summary CSV records whether the run was a \code{double} or
\code{float} build so the two can never be accidentally mixed.

\section{Unit tests}

\code{tests/unit\_tests.cpp} is a dependency-free test binary that is built in both the double and
the float configurations and run by \code{run\_all.sh} before any benchmark. It is small and
deliberately checks \emph{invariants} rather than reproducing expected numbers:

\begin{itemize}[leftmargin=*,topsep=2pt,itemsep=2pt]
\item \textbf{Linear algebra.} $\mA\mA^{-1}=\mI$ for a symmetric positive-definite $3\times3$
matrix; $\mL\mL^\T = \mA$ for the Cholesky factor; $\mR^\T\mR = \mB^\T\mB$ for the Householder QR of
a $5\times3$ matrix (the property the array square-root filters rely on); a rank-one Cholesky update
followed by the matching downdate returns the original factor; \code{solve} and \code{solve\_spd}
agree.
\item \textbf{Pole placement and Riccati.} Ackermann's formula places the poles of $(\mA-\mL\mC)$
where asked, verified through Cayley--Hamilton, for both a well-conditioned $3\times3$ pair and a
$4\times4$ battery-like pair with all poles near $z=1$ (the case the conditioned variant of
\S\ref{sec:lib-linalg} exists for); the doubling solution satisfies the discrete algebraic Riccati
equation to $10^{-8}$ and agrees with the fixed-point iteration on both pairs.
\item \textbf{Model Jacobians and OCV.} The analytic $\mF$, $\mH$, $\mB$ agree with central
differences to $10^{-6}$--$10^{-5}$ (double build only); the closed-form OCV is strictly monotone,
physically ranged at both ends, agrees with the 241-point table to \SI{3}{\milli\volt} and with its
analytic slope; the SPM stoichiometry map is consistent with the nominal capacity through the
lithium balance.
\item \textbf{Plants.} A 1\,C discharge for \SI{30}{\minute} removes exactly \num{0.5} of SOC in
both plants; after a rest the terminal voltage returns to the equilibrium OCV; the SPM voltage stays
within \SIrange{3.0}{4.25}{\volt}; the Preisach operator saturates symmetrically at $\pm M_p$, with
its descending branch above its ascending branch in the raw relay output --- which the plant negates,
so that the \emph{charging} branch of the hysteresis voltage lies above the \emph{discharging}
branch, as in a real cell.
\item \textbf{End-to-end sanity.} The EKF and UKF on the \code{init\_error} dataset must not
diverge, must reach $\text{RMSE}_{\text{ss}} < 0.02$ and must converge within \SI{900}{\second}.
\end{itemize}

The tests guard against regression; they are not a proof of correctness, and only the last touches
an estimator. Correctness of the estimators themselves is addressed by the cross-validation that
follows.

\section{Cross-validation against independent implementations}
\label{sec:me-crossval}

A library that implements seventy estimators against its own model, measured by its own metrics,
can be internally consistent and uniformly wrong. The only remedy is to reproduce a subset of the
results with software written by other people. Three independent reference implementations are
used, covering the three distinct things that could be wrong: the \emph{filters}, the
\emph{plant}, and the \emph{attitude conventions}. The procedures are given in full in
Appendix~\ref{app:validation}; this section reports what was found. All three comparisons are
driven by \code{tools/export\_validation}, which writes the datasets and the C++ trajectories in
full \code{\%.17g} precision so that the comparison is not limited by the file format.

\subsection{FilterPy: the EKF, the UKF and the RTS smoother}

FilterPy~\cite{labbe2014filterpy,labbe2020book} is the reference open-source Kalman library of the
Python ecosystem. The 2-RC battery model of Chapter~\ref{ch:ecm} was re-implemented from scratch in
NumPy~\cite{harris2020numpy} --- same parameters, same 241-point OCV table read from the same
exported CSV, same Arrhenius scaling, same $\mQ$ and $\mR$ --- and FilterPy's
\code{ExtendedKalmanFilter}, \code{UnscentedKalmanFilter} and \code{rts\_smoother} were run on the
nominal eVTOL dataset (\num{20100} samples).

% The auto-generated FilterPy table contains the literal identifier rts_smoother
% with an unescaped underscore. Rather than hand-edit a generated file, it is read
% with `_' made active so that it typesets as an underscore in text mode and still
% acts as a subscript in math mode. The file name is built first, while `_' still
% has an ordinary catcode, so that \input can scan it.
\begingroup\catcode`\_=12\relax
\gdef\estkitvalfilterpy{tables/validation_filterpy.tex}%
\endgroup
\begingroup\lccode`\~=`\_\relax
\lowercase{\endgroup\protected\gdef~{\ifmmode\sb\else\textunderscore\fi}}%
\begingroup
\catcode`\_=\active
\expandafter\input\expandafter{\estkitvalfilterpy}%
\endgroup

The EKF agrees to $10^{-15}$ in the SOC estimate and $10^{-19}$ in its variance --- round-off, over
\num{20100} steps, for two implementations that share no code. Reaching that level required
reproducing four ordinarily invisible conventions: the \emph{evaluation point of $\mQ$}
(\code{Ekf::predict} propagates the mean first and forms $\mQ(\xminus_k,\vu_{k-1})$ at the
\emph{predicted} state, whereas FilterPy adds a pre-set \code{self.Q} inside \code{predict()}, so
the driver evaluates the predicted mean itself and sets \code{ekf.Q} from it); \emph{symmetrisation}
of $\mP$ after both updates; the \emph{state constraint} $\soc\in[-0.05,1.05]$, $h\in[-1,1]$ applied
after the measurement update; and the \emph{input timing} convention of
\S\ref{sec:me-timing-convention}. Appendix~\ref{app:validation} gives each in full.

\paragraph{Two unscented variants, and why the difference matters.} The UKF agrees to
$2.5\times10^{-12}$ --- three orders of magnitude worse than the EKF, which is exactly what one
expects from a method whose update involves a Cholesky factorisation and a weighted sum over nine
sigma points, and which is still round-off. But the agreement is only obtainable because the
\code{estkit} UKF is configured for the comparison with
\code{opts.redraw\_sigma\_points = false}. The library supports two genuinely different unscented
updates:
\begin{itemize}[leftmargin=*,topsep=2pt,itemsep=2pt]
\item \textbf{Redraw} (the default, and what the benchmark uses): before the measurement update the
sigma points are regenerated from $(\xminus_k,\Pminus_k)$. The additive process noise $\mQ$, which
was added to $\mP$ after the propagation, is then reflected in the spread of the points and
therefore in $\mP_{yy}$. This is the form in Särkkä~\cite[Alg.~5.14]{sarkka2013} and in
Plett~\cite{plett2006spkf1}.
\item \textbf{Reuse} (\code{redraw\_sigma\_points = false}): the propagated points
$\mathcal{X}^f$ are reused for the measurement. This is Wan and van der
Merwe's original formulation~\cite[Alg.~3.1]{wan2000ukf} and is what FilterPy implements.
\end{itemize}
The two are not algebraically equivalent: reuse under-represents $\mQ$ in $\mP_{yy}$ and is
therefore slightly over-confident, while redrawing costs one extra Cholesky factorisation per step.
Reporting a single ``UKF'' number without saying which one is being used is a common and avoidable
source of irreproducibility in the literature; Chapter~\ref{ch:ukf} quantifies the difference on
this benchmark. A second, subtler point had to be handled at $k=0$: FilterPy's \code{update()}
reads \code{self.sigmas\_f}, which is still zero before the first \code{predict()}, so the driver
explicitly seeds it from $(\vx_0,\mP_0)$ --- which is what \code{estkit} does internally when
nothing has been propagated yet.

\paragraph{The RTS smoother and an affine change of variables.} FilterPy's \code{rts\_smoother}
implements the strictly \emph{linear} Rauch--Tung--Striebel recursion~\cite{rauch1965rts}: it
assumes $\vx_{k+1} = \mF_k\vx_k + \vw_k$ with no input term, and reconstructs the predicted mean
internally as $\mF_k\xplus_k$. The battery model has an input, but it is \emph{affine} in the
state:
\begin{equation}
\vx_{k+1} \;=\; \mF_k \vx_k + \bm{b}_k + \vw_k, \qquad
\mF_k = \diag\bigl(1,\,a_1(T_k),\,a_2(T_k),\,a_h(i_k)\bigr),\quad
\bm{b}_k = f(\bm 0,\vu_k),
\label{eq:me-affine}
\end{equation}
with $\mF_k$ independent of $\vx$. Define the deterministic particular solution
\begin{equation}
\bm{c}_0 = \bm 0, \qquad \bm{c}_{k+1} = \mF_k\bm{c}_k + \bm{b}_k,
\label{eq:me-particular}
\end{equation}
and the shifted state $\tilde{\vx}_k = \vx_k - \bm{c}_k$. Then
$\tilde{\vx}_{k+1} = \mF_k\tilde{\vx}_k + \vw_k$ is purely linear with the \emph{same} $\mF_k$ and
$\mQ_k$, and the covariances are unchanged because the shift is deterministic. Running the linear
smoother on $\{\xplus_k - \bm{c}_k\}$ and adding $\bm{c}_k$ back recovers the affine smoother exactly.
The residual disagreement, $7.3\times10^{-9}$, is three orders larger than the UKF's, for two
reasons that are worth naming: the backward pass forms differences of nearly equal $O(1)$
quantities after $\bm{c}_k$ has grown to $O(1)$, so there is genuine cancellation; and the forward
pass applied the state clamp, which is a non-linearity that the change of variables does not
commute with. Both effects are far below the \num{1e-4} scale at which any SOC result is quoted.

\subsection{PyBaMM: the electrochemical plant}

The filters can be right and the plant wrong. PyBaMM~\cite{sulzer2021pybamm} is the reference
open-source battery-modelling framework; its \code{lithium\_ion.SPM} with the built-in
\code{Chen2020} parameter set~\cite{chen2020parameterization} is the same single-particle model
that \code{plant\_spm.hpp} implements (Chapter~\ref{ch:spm}). For the comparison the
\code{estkit} SPM is configured with its published-parameter settings
($r_{\text{scale}} = k_{\text{scale}} = 1$, zero lumped ohmic resistance, isothermal) and both are
discharged at constant current from \SI{100}{\percent} SOC at \SI{25}{\celsius}; the
comparison table is \cref{tab:val_pybamm} in Chapter~\ref{ch:spm}.

The two implementations agree to \SIrange{0.6}{3.2}{\milli\volt} RMSE over the full discharge and
reach the \SI{2.5}{\volt} cut-off at the \emph{same second} at all three C-rates, which is the
stronger statement: the capacity, the stoichiometry mapping and the lithium balance are identical.
The maximum differences (\SIrange{13}{46}{\milli\volt}) all occur in the last seconds before
cut-off, where $|\mathrm{d}v/\mathrm{d}t|$ exceeds \SI{50}{\milli\volt\per\second} and a
sub-second timing difference becomes a large voltage difference; away from the knee the difference
is under \SI{5}{\milli\volt} everywhere. The residual is discretisation, not physics: \code{estkit}
uses 20 finite-volume shells per particle with backward Euler at $\dt=\SI{0.1}{\second}$ and a
linear surface-concentration reconstruction from the flux boundary condition, while PyBaMM uses its
default mesh and an adaptive DAE solver~\cite{andersson2019casadi}. The error grows with C-rate
exactly as a fixed-mesh solid-diffusion discretisation should, because the concentration gradient
in the particle steepens.

\subsection{The \texttt{ahrs} library: Madgwick, Mahony and the frame conventions}

The third comparison targets the part of the library where a mistake is most likely to be silent:
attitude conventions. The \code{ahrs} Python library~\cite{ahrs2019} implements the published
Madgwick~\cite{madgwick2011} and Mahony~\cite{mahony2008} algorithms. Both were run on the exported
nominal flight dataset with the same gains as the \code{estkit} implementations ($\beta=0.1$;
$k_P=1$, $k_I=0.1$), with accelerometer gating disabled on the \code{estkit} side so that the two
implement the same algorithm.

% Copyright 2026 Sreeram Anil
% SPDX-License-Identifier: CC-BY-4.0
\begin{table}[H]\centering\footnotesize\setlength{\tabcolsep}{4pt}\caption{Cross-validation of the estkit Madgwick ($\beta=0.1$) and Mahony ($k_P=1$, $k_I=0.1$) filters against the \texttt{ahrs} 0.4.0 Python library on the nominal flight dataset (90000 samples, \SI{100}{\hertz}): angular distance between the two implementations' quaternions, and attitude RMS error of each against the truth.}\label{tab:val_ahrs}
\begin{tabular}{lcccc}\toprule filter & max $\angle(\hat q_{\mathrm{estkit}}, \hat q_{\mathrm{ahrs}})$ [deg] & mean [deg] & RMS err.\ estkit [deg] & RMS err.\ ahrs [deg] \\ \midrule
Madgwick & 0.0014 & 0.0004 & 17.604 & 17.604 \\
Mahony & 0.0014 & 0.0004 & 22.639 & 22.639 \\
\bottomrule\end{tabular}\end{table}


The two implementations agree to \SI{0.0014}{\degree} peak and \SI{0.0004}{\degree} mean over
\num{90000} samples, and --- more tellingly --- produce attitude RMS errors against the truth that
are identical to three decimal places (\SI{17.604}{\degree} for Madgwick, \SI{22.639}{\degree} for
Mahony). Those large errors are \emph{not} implementation defects. They are the coordinated turns
of \S\ref{sec:ds-imu}: an ungated gravity-referenced filter reads a \SI{0.47}{}g sustained lateral
specific force as tilt and rolls its estimate into it. Reproducing the same large error with
independent code is a much stronger validation than reproducing a small one, because it confirms
that the dataset, the conventions and the algorithm all match; the gated variants registered in the
benchmark (Chapter~\ref{ch:att-madgwick}, Chapter~\ref{ch:att-mahony}) do far better, and the
difference is attributable entirely to the gate.

Two convention mismatches had to be resolved, and both are traps worth recording.

\emph{Specific force versus gravity.} \code{estkit} stores the accelerometer channel as specific
force, $f = R^\T(a - g_{\text{nav}})$, which reads $[0,0,-g]^\T$ in the FRD body frame at rest.
The \code{ahrs} filters expect an accelerometer that reads $+g$ along body $z$ at rest. The driver
therefore negates the channel before passing it to \code{ahrs}. A sign error here does not crash;
it produces a filter that converges smoothly to an upside-down attitude.

\emph{\code{ahrs}'s Mahony is hard-coded ENU.} While \code{ahrs.filters.Madgwick} builds its
magnetic reference from the measurement, \code{ahrs.filters.Mahony} constructs it as
$b = [\,0,\ \sqrt{h_x^2+h_y^2},\ h_z\,]$ --- north along $+y$, i.e.\ an East--North--Up navigation
frame --- which is incompatible with the NED/FRD dataset. The two frame pairs are related by the
proper rotation
\begin{equation}
\mP_{\text{ENU}\leftarrow\text{NED}} =
\begin{bmatrix} 0&1&0\\ 1&0&0\\ 0&0&-1\end{bmatrix},
\qquad \det \mP = +1,
\label{eq:me-enu}
\end{equation}
a \SI{180}{\degree} rotation about the axis $(1,1,0)/\sqrt2$, whose quaternion is
$q_P = [\,0,\ \tfrac{1}{\sqrt2},\ \tfrac{1}{\sqrt2},\ 0\,]$. The same matrix maps FRD to FLU, so all
three sensor vectors transform as $v_{\text{ENU}} = \mP v_{\text{NED}}$ and the attitude as
$q_{\text{ENU}} = q_P \otimes q_{\text{NED}} \otimes q_P^{-1}$. The driver rotates the inputs and
the initial quaternion into ENU/FLU, runs \code{ahrs}, and rotates the output back. In the ENU/FLU
pair the specific force is already what \code{ahrs} expects, so the earlier negation is undone.

\emph{Propagation.} A last, irreducible difference: \code{estkit} propagates the quaternion with
the exact exponential map $\hat q_{k} = \hat q_{k-1}\otimes\exp(\bm\omega\dt)$, while \code{ahrs}
takes a first-order Euler step of $\dot q = \frac12 q\otimes[0,\bm\omega]$. The per-step discrepancy
is $O\bigl((\|\bm\omega\|\dt)^2\bigr)$; at the dataset's peak rate of \SI{14.3}{\degree\per\second}
and $\dt=\SI{0.01}{\second}$ that is about \SI{5e-5}{\degree} per step, which the filters' own
feedback keeps bounded at the observed \SI{0.0014}{\degree}.

\subsection{What the cross-validation does and does not establish}

It establishes that the \code{estkit} implementations of the EKF, the UKF, the RTS smoother, the
SPM plant, Madgwick and Mahony compute what the published algorithms specify, to round-off, and
that the frame and timing conventions of the datasets are what this report says they are. By
extension it validates the shared infrastructure --- the linear algebra, the OCV table, the model
evaluation, the harness timing convention --- on which all seventy estimators depend, since an
error there would have shown up in the EKF comparison.

It does \emph{not} establish that the remaining sixty-odd estimators are correct; for those the
evidence is the unit tests, the internal consistency checks (the information filter, the
square-root EKF and the $U$--$D$ filter must reproduce the EKF to round-off, and do, which is
reported in Chapter~\ref{ch:ekf}), and the derivations in their own chapters. Nor does it establish
that the \emph{plants} are realistic: agreement with PyBaMM means two single-particle models agree,
not that a single-particle model is a lithium-ion cell. The limits of the modelling are discussed
in Chapters~\ref{ch:ecm} and~\ref{ch:spm}, and the standing recommendation of this report is that
any conclusion drawn from it be re-checked on measured data from the actual cell before it is used
in a design.

\section{Threats to validity}

Four are worth stating explicitly, because they bound how far the rankings in Parts~II--XIII can be
carried. \emph{Synthetic data}: every dataset is generated by a model, and for the ECM plant the
estimator-side model is a simplification of the same family. The SPM plant exists to break that
kinship --- a different model class, to which the estimators' ECM is fitted with a
\SI{4}{\milli\volt} residual (Chapter~\ref{ch:spm}) --- and the two plants are reported separately
everywhere for exactly this reason; rankings that hold on both are far more trustworthy than
rankings that hold on one. \emph{Tuning}: every estimator is run with the options its chapter
documents, taken from the literature's recommended values and from the sensor data sheet, never
optimised per scenario --- a choice for reproducibility and against peak performance, so a method
that needs careful per-scenario tuning looks worse here than in its original paper, which is the
intended signal for an implementer who has twelve cells and one set of gains. \emph{Scenario
coverage}: twelve cell, six pack and six attitude scenarios are a small sample of a large space, and
they are switched on one at a time (\code{combined} excepted), so interactions between faults are
largely untested. \emph{Single-run statistics}: three seeds is enough to see whether a difference is
robust and not enough to put a confidence interval on it, so differences below about
\SI{0.02}{\percent} SOC between two estimators should not be interpreted.
